% Near-verbatim transcript section: OY_napMPywE, [00:35:00.000,00:45:00.000)
% Canonical transcript: work/pilot/transcript.cleaned.jsonl
% Canonical transcript SHA-256: 5ac8ac5fb25a3235d8fa11b2b6be99b5f2bb9329d307c4045629544f4e43e9bd
% Exact notes: work/pilot/notes-exact.tex, original note pp. 3--4, combined PDF pp. 8--9
% Zero-word generated coverage records are omitted.

% YIN-VERBATIM: id=YIN-OY-T000115; video=OY_napMPywE; time=00:35:00.000-00:35:28.329; disposition=unresolved; notes=253a:3; pdf=8
\noindent\textbf{Room audio:} \textit{[Low-volume student exchange and pause; unintelligible.]}\par

% YIN-VERBATIM: id=YIN-OY-T000117; video=OY_napMPywE; time=00:35:28.339-00:35:48.250; disposition=included; notes=253a:3; pdf=8
\noindent\textbf{Yin:} Maybe a rather boring one, but the reason I'm talking about this is just to explain that you can easily construct a Lorentz-invariant quantum-mechanical system without ever talking about fields. Now, here it is.\par

% YIN-VERBATIM: id=YIN-OY-T000119; video=OY_napMPywE; time=00:35:48.260-00:36:08.210; disposition=included_with_uncertainty; notes=253a:3; pdf=8
\noindent\textbf{Yin:} There's a question about, you know, creation, how to create the [state]. I said, you know, that's not something I need to discuss, but we can, if you wish, define [creation and annihilation] operators in the following way at this point. So we can...\par

% YIN-VERBATIM: id=YIN-OY-T000121; video=OY_napMPywE; time=00:36:09.480-00:36:25.569; disposition=included_with_uncertainty; notes=253a:3; pdf=8
\noindent\textbf{Yin:} Introduce particle creation and [annihilation operators].\par

% YIN-VERBATIM: id=YIN-OY-T000123; video=OY_napMPywE; time=00:36:25.579-00:36:58.190; disposition=included; notes=253a:3; pdf=8; formula_class=EQUATION_NORMALIZED
% YIN-NOTE-EQUATION-PLACEHOLDER: notes=253a:3; pdf=8; exact_note={creation a_{\vec p}^{+}; annihilation a_{\vec p}}; transcript_spoken_form=dagger
\noindent\textbf{Yin:} Usually call these $a^\dagger$ and $a$, which are going to be also labeled by the momentum of the particle. The [label] I'm going to write as a subscript $\vec p$. Note that to specify the momentum of the particle I just need to specify the spatial momentum, and the energy is always, in this case, determined by the dispersion relation. So I don't have to say in addition what the energy is.\par

% YIN-VERBATIM: id=YIN-OY-T000125; video=OY_napMPywE; time=00:37:00.480-00:37:22.849; disposition=included; notes=253a:3; pdf=8; formula_class=EQUATION_NORMALIZED
% YIN-NOTE-EQUATION-PLACEHOLDER: notes=253a:3; pdf=8; exact_note={\ket{\vec p}=a_{\vec p}^{+}\ket{\Omega}}; transcript_spoken_form=dagger
\noindent\textbf{Yin:} So I'm going to define the creation operator by how it acts on these states. Well, the way that you can define that is more or less obvious. So, for example, the one-particle state should be the creation operator $a^\dagger_{\vec p}$ on the vacuum: $|\vec p\rangle=a^\dagger_{\vec p}|\Omega\rangle$.\par

% YIN-VERBATIM: id=YIN-OY-T000127; video=OY_napMPywE; time=00:37:22.859-00:37:38.510; disposition=included_with_uncertainty; notes=253a:3-4; pdf=8-9; formula_class=EQUATION_NORMALIZED
% YIN-NOTE-EQUATION-PLACEHOLDER: notes=253a:4; pdf=9; exact_note={\ket{\vec p_1,\vec p_2}=a_{\vec p_1}^{+}a_{\vec p_2}^{+}\ket{\Omega},\ \text{etc.}}; transcript_spoken_form=dagger
\noindent\textbf{Yin:} And if I would like to say, [for] a two-particle state, you just act two of these creation operators on the vacuum, and so forth: $|\vec p_1,\vec p_2\rangle=a^\dagger_{\vec p_1}a^\dagger_{\vec p_2}|\Omega\rangle$.\par

% YIN-VERBATIM: id=YIN-OY-T000129; video=OY_napMPywE; time=00:37:38.520-00:37:56.810; disposition=included; notes=253a:4; pdf=9; formula_class=EQUATION_NORMALIZED
% YIN-NOTE-EQUATION-PLACEHOLDER: notes=253a:4; pdf=9; exact_note_same_type_commutators_follow_in_YIN-OY-T000131
\noindent\textbf{Yin:} Now, I haven't so far specified the normalization of this basis of multiparticle states. Let me specify that now. It's important for me to specify the commutation relation of the $a$ and $a^\dagger$. By the way, of course, I should have said that these [commute].\par

% YIN-VERBATIM: id=YIN-OY-T000131; video=OY_napMPywE; time=00:38:00.260-00:38:19.849; disposition=included; notes=253a:4; pdf=9; formula_class=EQUATION_NORMALIZED
% YIN-NOTE-EQUATION-PLACEHOLDER: notes=253a:4; pdf=9; exact_note={[a,a]=[a^{+},a^{+}]=0}; transcript_spoken_form=dagger
\noindent\textbf{Yin:} And then I have to tell you the commutation relation of these creation and annihilation operators. For example, $[a_{\vec p},a_{\vec p'}]=[a^\dagger_{\vec p},a^\dagger_{\vec p'}]=0$.\par

% YIN-VERBATIM: id=YIN-OY-T000133; video=OY_napMPywE; time=00:38:22.140-00:38:51.890; disposition=included; notes=253a:4; pdf=9; formula_class=EQUATION_NORMALIZED
% YIN-NOTE-EQUATION-PLACEHOLDER: notes=253a:4; pdf=9; exact_note={[a_{\vec p},a_{\vec p'}^{+}]=\delta^{(D-1)}(\vec p-\vec p')}; transcript_spoken_form=dagger
\noindent\textbf{Yin:} There should be a nontrivial commutation relation between the annihilation operator and the creation operator. So we do $a_{\vec p}$ and $a^\dagger_{\vec p'}$, and this should be nonzero only if the particle momenta coincide. And so here let me choose the normalization convention. This is the delta function, so delta in $D-1$ spatial dimensions: $[a_{\vec p},a^\dagger_{\vec p'}]=\delta^{(D-1)}(\vec p-\vec p')$.\par

% YIN-VERBATIM: id=YIN-OY-T000135; video=OY_napMPywE; time=00:38:56.240-00:39:18.050; disposition=included_with_uncertainty; notes=253a:4; pdf=9; formula_class=EQUATION_NORMALIZED
% YIN-NOTE-EQUATION-PLACEHOLDER: notes=253a:4; pdf=9; exact_note_inner_product_follows_in_YIN-OY-T000139; creator_exact_note=a_{\vec p'}^{+}; transcript_spoken_form=dagger
\noindent\textbf{Yin:} As you can see, if I basically [take the] overlap, this is $\langle\vec p|\vec p'\rangle$, the one-particle state with momentum $\vec p'$, that I can write as... So this guy here is an $a^\dagger_{\vec p'}$ on the vacuum, and that guy...\par

% YIN-VERBATIM: id=YIN-OY-T000137; video=OY_napMPywE; time=00:39:18.060-00:39:40.750; disposition=included; notes=253a:4; pdf=9
% YIN-NOTE-EQUATION-PLACEHOLDER: notes=253a:4; pdf=9; exact_note={\braket{\vec p}{\vec p'}=\delta^{(D-1)}(\vec p-\vec p')}
\noindent\textbf{Yin:} There is a vacuum bra, actually, with the Hermitian conjugate, which is an $a_{\vec p}$. And then I use the commutation relation. I have to use this $a$, and like that I can replace [the product by the commutator]. So the result of this is $\delta^{(D-1)}(\vec p-\vec p')$.\par

% YIN-VERBATIM: id=YIN-OY-T000139; video=OY_napMPywE; time=00:39:40.760-00:39:53.930; disposition=included; notes=253a:4; pdf=9
% YIN-NOTE-EQUATION-PLACEHOLDER: notes=253a:4; pdf=9; exact_note={\braket{\vec p}{\vec p'}=\delta^{(D-1)}(\vec p-\vec p')}
\noindent\textbf{Yin:} So I'm just saying that I'm choosing a normalization for my one-particle state such that the overlap between one one-particle state and another one-particle state is this: $\langle\vec p|\vec p'\rangle=\delta^{(D-1)}(\vec p-\vec p')$.\par

% YIN-VERBATIM: id=YIN-OY-T000141; video=OY_napMPywE; time=00:39:53.940-00:40:20.510; disposition=included_with_uncertainty; notes=253a:4; pdf=9; formula_class=SOURCE_CONFLICT
% YIN-NOTE-EQUATION-PLACEHOLDER: notes=253a:4; pdf=9; exact_note_has_delta_only; spoken_extra_factor=[unclear-factor]; retain_conflict
\noindent\textbf{Student:} So why is it that sometimes, like, in some places I've seen this delta function has [unclear-factor]?\par
% YIN-VERBATIM: id=YIN-OY-T000141; video=OY_napMPywE; time=00:39:53.940-00:40:20.510; speaker=Yin; disposition=included_with_uncertainty; notes=253a:4; pdf=9; formula_class=SOURCE_CONFLICT
\noindent\textbf{Yin:} That is a matter of normalization convention, right? I'm free to rescale $a$ and $a^\dagger$ by a factor of two, and I would just change my normalization convention. Now, you see...\par

% YIN-VERBATIM: id=YIN-OY-T000143; video=OY_napMPywE; time=00:40:20.520-00:40:35.329; disposition=included; notes=253a:4; pdf=9
\noindent\textbf{Yin:} This one-particle state corresponds to a plane wave of the particle, which, strictly speaking, is not a normalizable state anyway. If you want to form a normalizable state, you have to take a wave packet. For example, you can take some state like that, you know, integrated by this...\par

% YIN-VERBATIM: id=YIN-OY-T000145; video=OY_napMPywE; time=00:40:35.339-00:41:18.890; speaker=Yin; disposition=included_with_uncertainty; notes=253a:4; pdf=9
\noindent\textbf{Yin:} ...against a function.\par
% YIN-VERBATIM: id=YIN-OY-T000145; video=OY_napMPywE; time=00:40:35.339-00:41:18.890; speaker=Student; disposition=included_with_uncertainty; notes=253a:4; pdf=9
\noindent\textbf{Student:} \textit{[Student question inaudible.]}\par
% YIN-VERBATIM: id=YIN-OY-T000145; video=OY_napMPywE; time=00:40:35.339-00:41:18.890; speaker=Yin; disposition=included_with_uncertainty; notes=253a:4; pdf=9
\noindent\textbf{Yin:} The overall [factor can be constant], or, you know, nonconstant. The overall factor is another convention. I'm allowed to, for example, also rescale this state by a factor of the energy to the one-third power, or whatever. It does not \textit{[remainder unintelligible].}\par

% YIN-VERBATIM: id=YIN-OY-T000147; video=OY_napMPywE; time=00:41:18.900-00:41:30.710; disposition=classroom_q_and_a; notes=253a:4; pdf=9
\noindent\textbf{Yin:} Any other questions?\par

% YIN-VERBATIM: id=YIN-OY-T000149; video=OY_napMPywE; time=00:41:34.160-00:41:48.829; speaker=Student; disposition=included_with_uncertainty; notes=253a:3-4; pdf=8-9
\noindent\textbf{Student:} Each state?\par
% YIN-VERBATIM: id=YIN-OY-T000149; video=OY_napMPywE; time=00:41:34.160-00:41:48.829; speaker=Yin; disposition=included_with_uncertainty; notes=253a:3-4; pdf=8-9
\noindent\textbf{Yin:} No, the annihilation operator should reduce the particle number by one, just like a creation operator should increase the particle number by one. So the vacuum has zero particles, and that's sort of the definition, part of the definition.\par

% YIN-VERBATIM: id=YIN-OY-T000151; video=OY_napMPywE; time=00:41:48.839-00:42:24.050; speaker=Yin; disposition=included; notes=253a:3; pdf=8
\noindent\textbf{Yin:} Any other questions?\par
% YIN-VERBATIM: id=YIN-OY-T000151; video=OY_napMPywE; time=00:41:48.839-00:42:24.050; speaker=Student; disposition=included; notes=253a:3; pdf=8
\noindent\textbf{Student:} Yes, but we also agreed to choose the energy of the vacuum state however we want, or...\par
% YIN-VERBATIM: id=YIN-OY-T000151; video=OY_napMPywE; time=00:41:48.839-00:42:24.050; speaker=Yin; disposition=included; notes=253a:3; pdf=8
% YIN-NOTE-EQUATION-PLACEHOLDER: notes=253a:3; pdf=8; exact_note={H\ket{\Omega}=0}
\noindent\textbf{Yin:} Right, so that's a good question. I stated this as following from our assumption that the vacuum is Poincare invariant. Now, of course, you could trivially redefine your Hamiltonian by adding to it a constant. You're allowed to do that. But I'm going to choose this to be zero. It's a natural assumption.\par

% YIN-VERBATIM: id=YIN-OY-T000153; video=OY_napMPywE; time=00:42:25.920-00:42:55.790; disposition=included_with_uncertainty; notes=253a:3; pdf=8
\noindent\textbf{Yin:} We would like to say that there is a [vacuum], but not even necessarily unique. I would assume that this vacuum state is symmetry-invariant, and we'll see shortly, in a bit, that this $H$ is one of the generators of our symmetry. And to say the state is invariant under the symmetry means all the generators should annihilate it, as in any continuous global symmetry.\par

% YIN-VERBATIM: id=YIN-OY-T000155; video=OY_napMPywE; time=00:42:55.800-00:43:00.910; disposition=classroom_q_and_a; notes=253a:3-4; pdf=8-9
\noindent\textbf{Yin:} Any other questions?\par

% YIN-VERBATIM: id=YIN-OY-T000157; video=OY_napMPywE; time=00:43:00.920-00:43:21.309; disposition=included; notes=253a:4; pdf=9
\noindent\textbf{Yin:} All right. So, given these creation and annihilation operators, well, why do I do this? Well, we can then express the Hamiltonian directly in terms of these creation and annihilation operators without having to say how it acts on the basis of states.\par

% YIN-VERBATIM: id=YIN-OY-T000159; video=OY_napMPywE; time=00:43:21.319-00:43:48.530; disposition=included; notes=253a:4; pdf=9; formula_class=EQUATION_NORMALIZED
% YIN-NOTE-EQUATION-PLACEHOLDER: notes=253a:4; pdf=9; exact_note={H=\int d^{D-1}\vec p\,\sqrt{\vec p^{\,2}+m^2}\,a_{\vec p}^{+}a_{\vec p}}; transcript_spoken_form=dagger
\noindent\textbf{Yin:} So you can easily verify that, in fact, in terms of $a$ and $a^\dagger$, we can write a Hamiltonian as an integral: $H=\int d^{D-1}\vec p\,\sqrt{\vec p^{\,2}+m^2}\,a^\dagger_{\vec p}a_{\vec p}$.\par

% YIN-VERBATIM: id=YIN-OY-T000161; video=OY_napMPywE; time=00:43:48.540-00:44:05.270; disposition=included_with_uncertainty; notes=253a:4; pdf=9; formula_class=EQUATION_NORMALIZED
% YIN-NOTE-EQUATION-PLACEHOLDER: notes=253a:4; pdf=9; exact_note_number_operator={a_{\vec p}^{+}a_{\vec p}}; transcript_spoken_form=dagger
\noindent\textbf{Yin:} The $a^\dagger_{\vec p}a_{\vec p}$ is the occupation number for the particles. So you can, I'll let you verify for yourself using the commutation relation, that this Hamiltonian [agrees with] the definition and produces the expected energy.\par

% YIN-VERBATIM: id=YIN-OY-T000163; video=OY_napMPywE; time=00:44:05.271-00:44:33.559; disposition=unresolved; notes=253a:4; pdf=9
\noindent\textbf{Student:} \textit{[Low-volume student question; unintelligible.]}\par

% YIN-VERBATIM: id=YIN-OY-T000165; video=OY_napMPywE; time=00:44:33.560-00:45:00.000; disposition=included_oral_qualification; notes=253a:4; pdf=9
\noindent\textbf{Yin:} That's much less obvious. We'll see that there are massless particles \textit{[unclear clause]}, and we can also study what we call a massless theory. And for massless particles, it's not obvious why you want to demand the number of particles to be finite, and that is a very subtle question. We may not even be able to get to it this semester, but it's definitely tied to the issue of infrared divergences in quantum electrodynamics.\par

% YIN-VERBATIM-COUNTS: eligible_words=1010; retained_words=1010; retention_ratio=1.000000
